% Zone „Das kennst du schon“ – aus bank/matrizen-und-uebergangsprozesse/zone.jsonl (zusammenbau v0.1)
\einheitenkopf[zone][Kennst du schon]{Das kennst du schon}
\setcounter{aufgabe}{0}

%% TODO Titel: Ich-kann-Satz fehlt in der Bank (Platzhalter: Kettenname) – Nr. 1
%% TODO Anweisung: ein Satz über der ersten Teilaufgabe fehlt in der Bank – Nr. 1
\begin{aufgabe}{Lineare Gleichungssysteme mit zwei und drei Unbekannten lösen, auch unterbestimmte (Parameter frei wählen) und überbestimmte (überzählige Gleichung als Probe)}
\begin{gleichungsraster}[2]
\gl{\text{Löse das Gleichungssystem: $x + y = 7$ und $x - y = 3$.}} &
\gl{\text{Löse das Gleichungssystem: $0{,}4x + 0{,}5y = 7$ und $x + y = 16$.}} \\
\end{gleichungsraster}
\end{aufgabe}

%% TODO Titel: Ich-kann-Satz fehlt in der Bank (Platzhalter: Kettenname) – Nr. 2
%% TODO Anweisung: ein Satz über der ersten Teilaufgabe fehlt in der Bank – Nr. 2
\begin{aufgabe}{Vektoren als Listen lesen und schreiben (Zustands-, Mengen-, Kostenvektoren), Komponenten adressieren}
\begin{teile}
\teil Ein Kiosk verkauft am Montag 40 Brezeln, 25 Brötchen und 12 Croissants. Schreibe die Mengen als Vektor (Brezeln $\mid$ Brötchen $\mid$ Croissants). ( \leerfeld | \leerfeld | \leerfeld )
\teil Berechne $2 \cdot (1{,}5 \mid 0{,}4 \mid 3) + (0{,}5 \mid 1{,}2 \mid 1)$. ( \leerfeld | \leerfeld | \leerfeld )
\end{teile}
\end{aufgabe}

%% TODO Titel: Ich-kann-Satz fehlt in der Bank (Platzhalter: Kettenname) – Nr. 3
%% TODO Anweisung: ein Satz über der ersten Teilaufgabe fehlt in der Bank – Nr. 3
\begin{aufgabe}{Diagramme lesen und zeichnen (Knoten, Pfeile, Punktdiagramme, Kurvenformen unterscheiden)}
\begin{teile}
\teil Lies ab: Wie viele Gäste kamen am Dienstag?

\saeulenab[ymax=100,ystep=20,yfein=5,ylabel=Gäste]{Mo/35,Di/60,Mi/45}
\leerfeld Gäste
\teil Lies ab und entscheide: Um welchen Faktor wächst die Anzahl von Woche zu Woche? Ist das Wachstum linear oder exponentiell?

\liniendia[ymax=40,ystep=10,yfein=5,ylabel=Anzahl]{W1/5,W2/10,W3/20,W4/40}
Faktor \leerfeld
\end{teile}
\end{aufgabe}

%% TODO Titel: Ich-kann-Satz fehlt in der Bank (Platzhalter: Kettenname) – Nr. 4
%% TODO Anweisung: ein Satz über der ersten Teilaufgabe fehlt in der Bank – Nr. 4
\begin{aufgabe}{Terme mit einer Unbekannten aufstellen und gegen Schranken auswerten (linear, mit Randwerten)}
\begin{teile}
\teil Eine Taxifahrt kostet 4\,€ Grundpreis und 2\,€ je Kilometer. Wie viele Kilometer kannst du mit 20\,€ höchstens fahren? \leerfeld[km]
\teil Der Anteil $0{,}2 + 0{,}5p$ hängt von p ab, p liegt zwischen 0 und 1. Zwischen welchen Werten liegt der Anteil? Setze die Randwerte ein. zwischen \leerfeld und \leerfeld
\end{teile}
\end{aufgabe}

%% TODO Titel: Ich-kann-Satz fehlt in der Bank (Platzhalter: Kettenname) – Nr. 5
%% TODO Anweisung: ein Satz über der ersten Teilaufgabe fehlt in der Bank – Nr. 5
\begin{aufgabe}{Anteile, Prozentsätze und Anzahlen ineinander umrechnen (Bleibe- und Wechselanteile, Gegenanteil)}
\begin{teile}
\teil Von 200 Kunden wechseln 10\,\% zu einem anderen Geschäft. Wie viele wechseln? \leerfeld Kunden
\teil Von 450 Fahrgästen bleiben 84\,\% im Bus sitzen. Wie viele bleiben sitzen? \leerfeld Fahrgäste
\end{teile}
\end{aufgabe}

%% TODO Titel: Ich-kann-Satz fehlt in der Bank (Platzhalter: Kettenname) – Nr. 6
%% TODO Anweisung: ein Satz über der ersten Teilaufgabe fehlt in der Bank – Nr. 6
\begin{aufgabe}{Potenzen mit Basis unter eins fallen, Exponentialungleichungen durch Probieren oder Logarithmieren lösen, Ergebnisse sinnvoll aufrunden}
\begin{teile}
\teil Berechne $0{,}5^3$. \leerfeld
\teil Für welche kleinste natürliche Zahl n gilt $1{,}2^n > 2$? Probiere. n = \leerfeld
\end{teile}
\end{aufgabe}

%% TODO Titel: Ich-kann-Satz fehlt in der Bank (Platzhalter: Kettenname) – Nr. 7
%% TODO Anweisung: ein Satz über der ersten Teilaufgabe fehlt in der Bank – Nr. 7
\begin{aufgabe}{Lineare Gleichungssysteme mit zwei und drei Unbekannten lösen, auch unterbestimmte (Parameter frei wählen) und überbestimmte (überzählige Gleichung als Probe) – Fehler finden}
\begin{teile}
\teil Tom soll alle Lösungen von $3x + y = 12$ angeben. Er schreibt: „$x = 4$, $y = 0$ – das ist die Lösung.“ Finde den Fehler und gib alle Lösungen an.
\end{teile}
\end{aufgabe}

%% TODO Titel: Ich-kann-Satz fehlt in der Bank (Platzhalter: Kettenname) – Nr. 8
%% TODO Anweisung: ein Satz über der ersten Teilaufgabe fehlt in der Bank – Nr. 8
\begin{aufgabe}{Lineare Gleichungssysteme mit zwei und drei Unbekannten lösen, auch unterbestimmte (Parameter frei wählen) und überbestimmte (überzählige Gleichung als Probe) – selbst rechnen}
\begin{gleichungsraster}[2]
\gl{\text{Gib alle Lösungen von $2x + y = 10$ an (setze $x = t$). Nenne die Lösung für $t = 3$.}} & \\
\end{gleichungsraster}
\end{aufgabe}
